\chapter{Esencialidad}

La esencialidad genética puede entenderse como una propiedad dependiente del contexto genético y ambiental del organismo en estudio. Esta propiedad puede ser entendida como binaria o continua. Debido a los datos disponibles para \textit{Acinetobacter baumannii}, la esencialidad se abordará mediante una clasificación binaria, ya que los ensayos aportan una clasificación binaria \cite{EsencialidadRancati2018}.

\comentario{Agregar algo acá o borrar la ecuación}
\[
E(g,c) =
\begin{cases}
1 & \text{si la pérdida de función de } g \text{ impide la viabilidad en } c, \\
0 & \text{si la pérdida de función de } g \text{ no impide la viabilidad en } c.
\end{cases}
\]

Donde \textit{c} representa la combinación de organismo/cepa, condiciones experimentales y estado fisiológico.

\comentario{Hablar de los diferentes experimentos, como fueron hechos, que aporta cada uno.}
